% =====================================================================
% SYNTHETIC SAMPLE DRAFT — for demonstrating the academic-paper-revision skill.
% This is a fictional, made-up paper. Authors, institution, dataset, and all
% numbers are invented and carry no IP. It intentionally contains the classic
% "student-draft" problems the skill is meant to fix:
%   - RQ1/RQ2/RQ3 scaffolding instead of prose contributions
%   - generic section headings ("Approach", "Experiment 1")
%   - chronological (dated-list) related work
%   - zero equations though the scoring rule is central
%   - a \ref to a figure that has no \label (renders as ??)
%   - loose, hedgy prose that runs over the page limit
% Target venue for the demo: a generic ACL-family conference.
% =====================================================================
\documentclass[11pt]{article}

\usepackage[utf8]{inputenc}
\usepackage{graphicx}
\usepackage{booktabs}
\usepackage{url}

\title{Punctuation-Aware Tokenization for Sentiment Classification of Product Reviews}
\author{Jane Doe \and John Roe \\ Department of Sample Studies, University of Nowhere \\ \texttt{\{jdoe, jroe\}@example.edu}}
\date{}

\begin{document}
\maketitle

\begin{abstract}
Sentiment classification is a very important task in natural language
processing and a lot of people have worked on it. In this paper we look at
whether keeping punctuation during tokenization helps a simple classifier do
better on product reviews. We built a small pipeline and ran some experiments
on a dataset of reviews that we collected. Our results show that keeping
punctuation seems to help a little bit, and we think this is interesting and
could be useful for future work. We also tried a few different classifiers and
report what we found.
\end{abstract}

\section{Introduction}

Sentiment analysis is the task of deciding whether a piece of text is positive
or negative. It is used in many places like product reviews, social media, and
customer feedback. Many researchers have studied this problem over the years
and there are many methods. Most methods start by tokenizing the text, which
means splitting it into words or tokens. Usually punctuation is thrown away
during tokenization because it is thought to be noise.

In this paper we ask whether that is a good idea. Maybe punctuation carries
sentiment information. For example, an exclamation mark might mean the reviewer
is excited, and lots of question marks might mean they are confused or annoyed.
So we decided to study this.

We investigate the following research questions:
\begin{itemize}
    \item \textbf{RQ1:} Does keeping punctuation during tokenization improve
    sentiment classification accuracy?
    \item \textbf{RQ2:} Which punctuation marks are the most useful for the
    classifier?
    \item \textbf{RQ3:} Does the effect depend on which classifier we use?
\end{itemize}

To answer these questions we built a pipeline (see Figure~\ref{fig:pipeline})
and ran experiments. The rest of the paper is organized as follows. Section 2
talks about related work. Section 3 describes our approach. Section 4 and
Section 5 present our experiments. Section 6 concludes.

\section{Related Work}

In 2013, Mikolov et al. introduced word2vec which is a way to represent words
as vectors. In 2014, Pennington et al. introduced GloVe which is another way to
get word vectors. In 2015, some people used LSTMs for sentiment analysis and got
good results. In 2018, Devlin et al. introduced BERT which is a transformer
model and it became very popular for many NLP tasks including sentiment
analysis. After that, in 2019 and 2020, many people fine-tuned BERT and other
transformers for sentiment classification and kept improving the numbers.
There has also been work on tokenization, for example byte-pair encoding and
WordPiece, but most of this work is about handling rare words and not about
punctuation. So our work is different because we focus on punctuation.

\section{Approach}

Our approach has a few steps. First we take the raw review text. Then we
tokenize it. We have two settings: one where we remove punctuation (the
baseline) and one where we keep punctuation as separate tokens (our method).
Then we turn the tokens into features. We use a bag-of-words representation
where each feature is the count of a token. Then we train a classifier on these
features. We compute a score for each review and if the score is above a
threshold we call it positive, otherwise negative. We picked the threshold by
trying a few values and choosing the one that worked best on the training set.

We tried three classifiers: logistic regression, a naive Bayes classifier, and
a small feed-forward neural network. We used default settings for most things.

\begin{figure}[t]
    \centering
    \includegraphics[width=0.7\linewidth]{example-image}
    \caption{The pipeline of our approach. Raw text is tokenized, turned into
    bag-of-words features, and classified.}
\end{figure}

\section{Experiment 1}

We collected a dataset of 5,000 product reviews from an online store. Each
review has a star rating from 1 to 5. We treated 4 and 5 stars as positive and
1 and 2 stars as negative, and we threw away 3-star reviews because they are
ambiguous. This gave us about 4,200 reviews. We split them 80/20 into training
and test.

We ran the baseline (no punctuation) and our method (with punctuation) using
logistic regression. The results are in Table~\ref{tab:main}. Keeping
punctuation improved accuracy from 0.812 to 0.831. We think this shows that
punctuation is helpful.

\begin{table}[t]
    \centering
    \begin{tabular}{lc}
        \toprule
        Setting & Accuracy \\
        \midrule
        Baseline (no punctuation) & 0.812 \\
        Ours (with punctuation) & 0.831 \\
        \bottomrule
    \end{tabular}
    \caption{Main results with logistic regression.}
    \label{tab:main}
\end{table}

\section{Experiment 2}

To answer RQ2 and RQ3 we did more experiments. For RQ2 we looked at which
punctuation marks helped the most. We found that exclamation marks and question
marks were the most useful, while commas and periods did not help much. For RQ3
we tried the other two classifiers. Naive Bayes went from 0.788 to 0.795 and the
neural network went from 0.822 to 0.834. So the effect seems to hold across
classifiers, although the size of the improvement is different.

\section{Conclusion}

In this paper we studied whether keeping punctuation during tokenization helps
sentiment classification. We found that it helps a little bit across three
classifiers, and that exclamation marks and question marks are the most useful.
In future work we would like to try more datasets, more languages, and bigger
models. We think punctuation-aware tokenization is a promising direction.

\begin{thebibliography}{9}
\bibitem{mikolov2013} Mikolov, T. et al. (2013). Efficient Estimation of Word Representations in Vector Space.
\bibitem{pennington2014} Pennington, J. et al. (2014). GloVe: Global Vectors for Word Representation.
\bibitem{devlin2018} Devlin, J. et al. (2018). BERT: Pre-training of Deep Bidirectional Transformers for Language Understanding.
\end{thebibliography}

\end{document}
